En este capítulo se presentan los objetivos y el catálogo de requisitos del sistema **UniHub**. Opcionalmente, se incluye el análisis de la brecha entre los requisitos planteados y la solución base seleccionada. 

\section{Objetivos de Negocio}
El objetivo de negocio principal es ofrecer un portal público centralizado de enseñanza superior en España, facilitando la transparencia, la comparativa de planes de estudio BOE y el acceso a la oferta universitaria pública y privada.

\section{Objetivos del Sistema}
1. Automatizar la extracción periódica de universidades y titulaciones vigentes (Grados, Másteres y Doctorados).
2. Homogenizar y almacenar relacionalmente los datos en PostgreSQL ordenando públicamente primero las universidades públicas y posteriormente las privadas.
3. Exponer una API REST documentada en FastAPI con control de permisos y métricas físicas cgroup.
4. Construir un portal web interactivo con estética corporativa de la Universidad de Cádiz (UCA), paginación configurable (5, 10, 20, 50, 100) y ocultación de códigos internos.
5. Proporcionar un buscador por geolocalización (cercanía en km con distancia integrada directamente en tarjetas).
6. Recolectar métricas de uso y clasificar las 6 universidades más visitadas en "Universidades Destacadas".
7. Incorporar el apartado "Sobre Nosotros" (TFG Alejandro Ramos Rodríguez, UCA).
8. Aislar el panel de control del Administrador de la Web tras la ruta manual \texttt{/admin}.
9. Desplegar la infraestructura mediante contenedores Docker orquestados con persistencia garantizada.

\section{Requisitos funcionales}
\begin{itemize}
    \item \textbf{RF-01 Descarga de Catálogo Oficial:} Descargar e inspeccionar archivos XLS de universidades y titulaciones.
    \item \textbf{RF-02 Filtrado Estricto de Vigencia y Clasificación de Doctorados:} Excluir titulaciones extinguidas, no vigentes o eliminadas; conservar únicamente las vigentes o autorizadas y clasificar Doctorados (RD 99/2011).
    \item \textbf{RF-03 Parseo Meticuloso de BOE:} Escanear todos los PDF candidatos del BOE para extraer asignaturas, módulos, materias, créditos ECTS, curso y cuatrimestre.
    \item \textbf{RF-04 Resiliencia y Notificación Automática:} Aplicar pausas de 5 min (10 fallos) y cortocircuito a los 15 min ante caídas de red, notificando via POST a la API REST al finalizar.
    \item \textbf{RF-05 Regla de No Sobrescritura Vacía:} Evitar que campos vacíos o indeterminados (\texttt{""}, \texttt{"undefined"}) sobrescriban información previa válida.
    \item \textbf{RF-06 API REST de Consulta con Ordenación Prioritaria:} Endpoints para listar universidades (públicas primero, privadas después) y titulaciones por nivel académico.
    \item \textbf{RF-07 API REST CRUD de Administración y Métricas cgroup:} Endpoints \texttt{POST}, \texttt{PUT} y \texttt{DELETE} para universidades y titulaciones, junto con métricas de uso de RAM/CPU por contenedor.
    \item \textbf{RF-08 Buscador, Filtros y Paginación Web:} Búsqueda global por nombre, tipo (Pública/Privada), CCAA y nivel (Grado, Máster, Doctorado), con paginación de 5, 10, 20, 50 y 100 resultados.
    \item \textbf{RF-09 Ocultación de Códigos Internos:} Ocultar códigos internos numéricos en la interfaz pública web.
    \item \textbf{RF-10 Visor Modal de Planes BOE:} Visualización en la web del desglose ECTS y enlace oficial al PDF del BOE.
    \item \textbf{RF-11 Búsqueda por Cercanía con Distancia Integrada:} Cálculo de distancias en km mediante la fórmula de Haversine mostrando la distancia dentro de la tarjeta del centro.
    \item \textbf{RF-12 Universidades Destacadas (Top 6):} Cálculo dinámico de las 6 universidades más consultadas en la plataforma.
    \item \textbf{RF-13 Sección "Sobre Nosotros":} Página informativa del TFG de Alejandro Ramos Rodríguez (Ingeniería Informática, UCA).
    \item \textbf{RF-14 Panel del Administrador en Ruta Manual \texttt{/admin}:} Acceso aislado sin botón visible para gestionar CRUD y monitorear la salud individual de contenedores.
\end{itemize}

\section{Requisitos no funcionales}
\begin{itemize}
    \item \textbf{RNF-01 Rendimiento:} Latencia promedio de respuesta de la API REST inferior a 100 ms.
    \item \textbf{RNF-02 Usabilidad e Identidad Visual:} Interfaz web profesional con paleta corporativa de la Universidad de Cádiz (UCA).
    \item \textbf{RNF-03 Seguridad:} Rol de solo lectura en base de datos (\texttt{rupt\_api\_user}) y acceso restringido por ruta manual al panel de administración.
    \item \textbf{RNF-04 Resiliencia a Fallos:} El crawler no se detendrá ante errores individuales de red o parseo de PDFs, registrándolos en \texttt{errores\_crawler.json} y aplicando el cortocircuito de 5m/15m.
    \item \textbf{RNF-05 Persistencia y Portabilidad:} Despliegue en 4 contenedores Docker independientes con persistencia garantizada en volumen nominado y montaje directo de host.
\end{itemize}

\section{Requisitos de información}
El sistema gestiona las siguientes entidades principales de información: \textit{Universidad}, \textit{Titulación}, \textit{PlanEstudios}, \textit{ResumenCréditos}, \textit{ElementoCurricular}, \textit{ErrorCrawler} y \textit{EstadísticaRendimiento}.

\section{Reglas de negocio}
\begin{itemize}
    \item \textbf{RN-01 Jerarquía por Real Decreto:} En presencia de múltiples planes de una misma titulación, priorizar RD 822/2021 > RD 1393/2007 > RD 56/2005.
    \item \textbf{RN-02 No Omisión Entre Universidades:} Misma titulación en distintas universidades debe ser tratada como titulación independiente.
    \item \textbf{RN-03 Comprobación Incremental Estricta:} No volver a descargar un PDF del BOE si la URL y la fecha coinciden con el registro de \texttt{checkpoint.json}.
    \item \textbf{RN-04 Ordenación Prioritaria:} Mostrar siempre las universidades públicas en primer lugar y posteriormente las privadas.
\end{itemize}

\section{Análisis GAP}
Dado que **UniHub** se ha construido a medida mediante una arquitectura en 4 fases que integra rastreo, API REST propia, frontend con diseño UCA y orquestación Docker, no se requirió la adopción directa de plataformas base de terceros.